\section{Recovering the field from simultaneous values}\label{sec:field}

The preceding section recovers the good rational subfields and the points of
their projective lines.  We now recover the field operations from this
information.  A single true divisor can detect the values of several good
functions at once.  The resulting correspondence of algebraic relations
will first identify the constant fields and then give an isomorphism of
the perfect closures.

For every good function \(t\in K\), choose a generator \(t'\in L\) of the good
rational subfield matched with \(k(t)\) by
Proposition~\ref{prop:good-matching}.  Its point correspondence, expressed in
these coordinates, is a bijection
\[
 B_t:\Pj^1(k)\longrightarrow\Pj^1(l).
\]
Changing \(t'\) by a projective linear transformation, we arrange
\(B_t(\infty)=\infty\).  These choices are made independently for the
different good functions.  At this stage the maps \(B_t\) are only
bijections of sets.

\subsection{Simultaneous values}

For a finite tuple \(\mathbf t=(t_1,\ldots,t_m)\) of good functions, let
\(X_{\mathbf t}\subset(\Pj^1_k)^m\) be the reduced closure of the image of
the rational map defined by the \(t_i\) on a model of \(K/k\).  Thus
\(X_{\mathbf t}\) is the irreducible variety of algebraic relations among
the \(t_i\).  Define \(Y_{\mathbf t}\subset(\Pj^1_l)^m\) in the same way
from the chosen \(t_i'\).

\begin{lemma}[Simultaneous values]\label{lem:simultaneous-values}
For every finite tuple \(\mathbf t\) of good functions there are dense open
subsets \(U\subset X_{\mathbf t}\) and \(U'\subset Y_{\mathbf t}\) such that
\begin{equation}\label{eq:field-simultaneous}
 \left(\prod_i B_{t_i}\right)(U(k))\subset Y_{\mathbf t}(l),
 \qquad
 \left(\prod_i B_{t_i}^{-1}\right)(U'(l))\subset X_{\mathbf t}(k).
\end{equation}
\end{lemma}

\begin{proof}
Choose a smooth model open on which all the \(t_i\) are regular and all the
differentials \(dt_i\) are nowhere zero.  Such an open exists.  In positive
characteristic, \(t_i\) cannot be a \(p\)-th power in \(K\), since its
\(p\)-th root would be algebraic over the relatively algebraically closed
subfield \(k(t_i)\).  As \(k\) is perfect, this gives \(dt_i\ne0\); in
characteristic zero the same conclusion is immediate.

At a closed point \(P\) of this open, write \(a_i=t_i(P)\).  The exceptional
divisor of the blowup at \(P\) defines a true divisorial valuation \(v\) with
\[
 \ord_v(t_i-a_i)=1\qquad\text{for every }i.
\]
Indeed \(dt_i(P)\ne0\) says that \(t_i-a_i\) has order one in the maximal
ideal of the regular local ring at \(P\).  Thus restriction of \(I_v\)
to each \(k(t_i)\) is the nonzero point line at \(a_i\).

By Proposition~\ref{prop:true-divisors}, the matched valuation on \(L\) is a
true divisor.  Restriction of characters commutes with the correspondence
of good subfields, so its center in the \(t_i'\)-line is \(B_{t_i}(a_i)\)
for every \(i\).  Specializing the target tuple along this valuation shows
that \((B_{t_i}(a_i))_i\) belongs to \(Y_{\mathbf t}(l)\): the generic tuple
lies in that closed subvariety of the product of projective lines, and so
does its specialization.

The image of the chosen model open in \(X_{\mathbf t}\) is constructible
and dense, hence contains a dense open subset \(U\).  Every \(k\)-point of
\(U\) has a closed preimage, since its nonempty fiber is of finite type
over the algebraically closed field \(k\).  This proves the first inclusion
in \eqref{eq:field-simultaneous}.  The same argument for \(\Theta^{-1}\)
proves the second.
\end{proof}

These inclusions provide information in the Zariski topology without
asserting that any \(B_t\) preserves that topology.  We next apply them to
the relation \(z=t+cs\).  This will force the coordinate bijections to have
a common field-theoretic form.

\subsection{Aligning the constant fields}

An element of \(l(X_1,\ldots,X_r)^i\) will be called a \emph{perfect rational
function}.  It has a well-defined value at a general tuple of \(l\)-points:
in positive characteristic one first raises it to a sufficiently large
\(p\)-power and then takes the unique corresponding root in \(l\).
Here and below, \emph{general} means belonging to a suitable dense open.

Let \(\mathcal H_l\) be the group of permutations of \(\Pj^1(l)\) generated
by \(\PGL_2(l)\), together with \(x\mapsto x^p\) when \(\charac l=p>0\).

\begin{lemma}\label{lem:cofinite-rational}
A perfect rational function in one variable that is injective on a cofinite
subset of \(\Pj^1(l)\) agrees there with a unique member of \(\mathcal H_l\).
In characteristic zero,
\(\mathcal H_l=\PGL_2(l)\); in characteristic \(p>0\), every member has a
unique expression
\[
 M\circ\Frob^n,\qquad M\in\PGL_2(l),\quad n\in\Z.
\]
In particular, \(\mathcal H_l\) embeds in the group of permutations of
\(\Pj^1(l)\) modulo agreement outside finite sets.
\end{lemma}

\begin{proof}
In characteristic zero, injectivity on a cofinite set implies that a
rational map has degree one.  In characteristic \(p\), choose \(N\ge0\)
such that the given function \(h\) satisfies \(h^{p^N}=r\in l(X)\).
Write \(r=q^{p^a}\), with \(a\ge0\) maximal and \(q\in l(X)\) separating.
The map \(q\) is also injective on a cofinite set, so its degree is one.
Indeed a separating map of degree greater than one has more than one
distinct point over a general value, and deleting finitely many points
cannot change this.  Thus \(q\) is projective linear, and
\(h=\Frob^{a-N}\circ q\) on points.

Conjugating a projective linear transformation by Frobenius raises its
coefficients to \(p\)-th powers.  This proves the asserted normal form.
No nonzero power of Frobenius is projective linear: for a positive exponent
the corresponding map has inseparable degree greater than one, and a
negative exponent reduces to this case by inversion.  The normal form is
therefore unique.  Finally, two perfect rational functions agreeing
cofinitely are equal, as is seen after clearing Frobenius powers and
comparing ordinary rational functions.
\end{proof}

Write \(T_a(x)=x+a\) and \(m_c(x)=cx\), fixing \(\infty\); these
transformations generate \(\Aff(k)\).  We compare bijections of projective
point sets modulo finite disagreement.  Composition is well defined with
this convention, because bijections carry finite sets to finite sets.
Membership in \(\mathcal H_l\) in the next lemma means membership after
this identification.  Lemma~\ref{lem:cofinite-rational} ensures that the
representing member of \(\mathcal H_l\) is unique.

\begin{lemma}[Transported addition]\label{lem:field-affine}
Let \(t,s\) be algebraically independent good functions such that
\(z_c=t+cs\) is good for every \(c\in k^\times\).
For every \(c\in k^\times\) there is a perfect rational function
\(R_c\in l(X,Y)^i\) such that
\begin{equation}\label{eq:field-sum}
 R_c(x,y)=B_{z_c}\bigl(B_t^{-1}(x)+cB_s^{-1}(y)\bigr)
 \qquad\text{for general }(x,y)\in l^2.
\end{equation}
The one-variable slices of these functions give
\begin{equation}\label{eq:field-affine}
 B_sB_t^{-1}\in\mathcal H_l,
 \qquad
 B_t\Aff(k)B_t^{-1}\subset\mathcal H_l.
\end{equation}
The second inclusion is an embedding of groups.
\end{lemma}

\begin{proof}
Consider the relation variety \(Y\) for \((t,s,z_c)\).  Its projection to
the first two factors is dominant.  Indeed \(k(t)\) and \(k(s)\) are
distinct curve subfields, so their matched subfields \(l(t')\) and
\(l(s')\) are distinct.  Any two distinct curve subfields have compositum
of transcendence degree two: if their compositum had transcendence degree
one, relative algebraic closedness would make them equal.

On the dense open of \(Y\) given by the inverse inclusion in
Lemma~\ref{lem:simultaneous-values}, the third coordinate is forced by the
first two to be
\[
 B_{z_c}\bigl(B_t^{-1}(x)+cB_s^{-1}(y)\bigr).
\]
We restrict to finite coordinates, which is permitted because all the
coordinate functions are nonconstant and the \(B\)'s fix \(\infty\).
If \(\dim Y=3\), then \(Y=(\Pj^1_l)^3\), and a dense open has infinitely
many third coordinates above a general pair \((x,y)\), a contradiction.
Consequently \(\dim Y=2\).

The projection \(Y\to(\Pj^1_l)^2\) is proper, dominant and generically
finite.  The complement of the open just used has dimension at most one,
so its image is a proper closed subset of the base.  After removing that
image and restricting to the locus of finite fibers, every closed fiber
has precisely one point.  Thus the finite extension \(l(Y)/l(x,y)\)
has separable degree one.  Indeed, in positive characteristic write the
minimal polynomial of the third coordinate as \(g(Z^{p^e})\), with \(g\)
separable; in characteristic zero take the minimal polynomial itself.
After shrinking the base to preserve the degree and nonzero discriminant
of the separable polynomial,
the separable degree counts the distinct points of a general fiber,
since taking \(p^e\)-th roots is a bijection on algebraically closed
points.  The extension is therefore purely
inseparable, and the third coordinate gives a perfect rational function
\(R_c\) satisfying \eqref{eq:field-sum}.

Fixing a general \(y\) in \eqref{eq:field-sum} gives a perfect rational
function of \(x\), agreeing cofinitely with a bijection.  A dense open in
\((\Pj^1_l)^2\) has cofinite fibers outside finitely many values of \(y\).
Since \(B_s\) is a bijection, Lemma~\ref{lem:cofinite-rational} therefore gives
\[
 D_b:=B_{z_c}T_{cb}B_t^{-1}\in\mathcal H_l
 \qquad\text{for cofinitely many }b\in k.
\]
For two such parameters,
\[
 D_b^{-1}D_{b'}=B_tT_{c(b'-b)}B_t^{-1}.
\]
Every element of \(k\) is a difference of two members of a cofinite subset:
for any prescribed difference, the two required cofinite sets intersect.
Hence all translations conjugated by \(B_t\) belong to \(\mathcal H_l\).
The products \(D_bD_{b'}^{-1}\) give the same assertion for \(B_{z_c}\),
and then \(D_b\) gives \(B_{z_c}B_t^{-1}\in\mathcal H_l\).

Fixing a general \(x\) instead yields
\[
 B_{z_c}T_a m_cB_s^{-1}\in\mathcal H_l
 \qquad\text{for cofinitely many }a\in k.
\]
Composing on the left successively with
\[
 B_{z_c}T_{-a}B_{z_c}^{-1}
 \quad\text{and}\quad
 B_tB_{z_c}^{-1},
\]
both already in \(\mathcal H_l\), gives
\[
 B_t m_cB_s^{-1}\in\mathcal H_l\qquad(c\ne0).
\]
At \(c=1\) this gives \(B_sB_t^{-1}\in\mathcal H_l\).  Composing again
gives all conjugated multiplications \(B_t m_cB_t^{-1}\).  This proves
\eqref{eq:field-affine}.  Finally, two distinct affine transformations
cannot agree outside a finite set; conjugating by a bijection preserves
this property.  Thus the resulting homomorphism is injective.
\end{proof}

We have obtained an affine group action by perfect rational transformations.
Its translation subgroup will identify the addition on the constants.
The two-variable function \(R_1\) is then needed to show that the resulting
embedding of constant fields is onto.

\begin{proposition}[Alignment of the constants]\label{prop:constant-alignment}
The fields \(k\) and \(l\) have the same characteristic.  There is a field
isomorphism \(\sigma:k\longrightarrow l\) such that, for every good function
\(t\), there is \(H_t\in\mathcal H_l\) with
\begin{equation}\label{eq:field-alignment}
 B_t=H_t\circ\sigma
 \qquad\text{outside a finite subset of }\Pj^1(k),
\end{equation}
where \(\sigma(\infty)=\infty\).
\end{proposition}

\begin{proof}
Choose a good function \(t\), and use Lemma~\ref{lem:common-partner} to
choose a good \(s\) satisfying the hypotheses of
Lemma~\ref{lem:field-affine}.  First consider the embedding
\(B_t\Aff(k)B_t^{-1}\subset\mathcal H_l\), with the finite-agreement
convention of that lemma.

\emph{The affine image is projective linear.}
In positive target characteristic the Frobenius exponent defines a
homomorphism \(\mathcal H_l\to\Z\).  Its restriction to the affine image is
zero: \(k^\times\) is divisible, while the additive group of \(k\) is
divisible in characteristic zero and torsion in positive characteristic.
Both groups therefore have zero image in \(\Z\), and they generate
\(\Aff(k)\).  In characteristic zero for \(l\), there is no exponent to
consider.  We have in either case an embedding into \(\PGL_2(l)\).

\emph{Translations yield a field embedding.}
Let \(U\) be the image of the translation subgroup.  This is an infinite
abelian subgroup of \(\PGL_2(l)\), normalized by the whole affine image.
It cannot contain a nonidentity semisimple element.  To see this, such an
element has two fixed points, and every element of \(U\) preserves their
pair.  The subgroup fixing both points has index at most two in \(U\);
it is infinite, hence contains an element of order greater than two.
The centralizer of that element is the torus fixing the two points, so
all of \(U\) lies in that torus.  Any transformation normalizing \(U\)
preserves the same pair and acts on \(U\) by either identity or inversion.
This contradicts the faithful conjugation action of \(k^\times\) on its
additive group in \(\Aff(k)\).

Every nonidentity element of \(U\) is consequently unipotent.  Conjugate
the image in \(\PGL_2(l)\) so that one such element is a translation with
unique fixed point \(\infty\).  Its centralizer consists of translations,
and hence so does \(U\).  Write this conjugation as postcomposition
\(B=MB_t\), with \(M\in\PGL_2(l)\).  We obtain an injective additive map
\(\chi:k\to l\) for which
\begin{equation}\label{eq:field-translation}
 BT_bB^{-1}=T_{\chi(b)}
 \qquad\text{modulo finite disagreement}.
\end{equation}
After rescaling the target coordinate, assume \(\chi(1)=1\).
The image of \(m_c\) normalizes \(U\), hence fixes \(\infty\), and is an
affine transformation of some slope \(\lambda_c\).
Conjugating translations gives
\[
 \chi(cb)=\lambda_c\chi(b).
\]
Setting \(b=1\) gives \(\lambda_c=\chi(c)\).  Thus \(\chi\) is
multiplicative as well as additive and is a field embedding.  In
particular, the characteristics of \(k\) and \(l\) agree.

\emph{The field embedding is onto.}
Set \(z=t+s\).  Lemma~\ref{lem:field-affine} gives members
\(M_s,M_z\in\mathcal H_l\) such that
\[
 B_s=M_sB,\qquad B_z=M_zB
 \qquad\text{modulo finite disagreement}.
\]
Using \(R_1\) from \eqref{eq:field-sum}, define the perfect rational function
\[
 F(X,Y)=M_z^{-1}\bigl(R_1(M^{-1}X,M_sY)\bigr).
\]
For cofinitely many \(b\in k\), and for each such \(b\) for cofinitely many
\(X\in l\), equations \eqref{eq:field-sum} and
\eqref{eq:field-translation} give
\begin{equation}\label{eq:field-slice}
 F(X,B(b))=B\bigl(B^{-1}(X)+b\bigr)=X+\chi(b).
\end{equation}
Here is the order of the exclusions.  First exclude the finitely many
values of \(b\) at which \(B(b)\) is infinite, the equality \(B_s=M_sB\)
fails, or the corresponding vertical fiber misses the open of
\eqref{eq:field-sum}.  For each remaining \(b\), that open excludes only
finitely many \(X\); the equality \(B_z=M_zB\) excludes finitely many more,
since its argument is a bijective function of \(X\).  Finally
\eqref{eq:field-translation} excludes a finite set of \(X\), which may
depend on \(b\).

For each fixed remaining \(b\), \eqref{eq:field-slice} is therefore an
identity of perfect rational functions in \(X\).
Choose \(X_0\in l\) such that \(F(X_0,Y)\) is defined and finite for general
\(Y\).  This choice is possible by expressing a power of \(F\) as a
quotient of polynomials and choosing \(X_0\) for which its denominator
does not vanish identically as a polynomial in \(Y\).
Define
\[
 h(Y)=F(X_0,Y)-X_0.
\]
For all but finitely many of the remaining \(b\), the denominator is
nonzero at \(Y=B(b)\); evaluating the rational identity then gives
\[
 \chi(b)=h(B(b)).
\]
In particular, this evaluation does not require \(X_0\) to avoid the
original exceptional sets for all \(b\) at once.

The injectivity of \(\chi\) implies that \(h\) is injective on a cofinite
subset of \(\Pj^1(l)\).  By Lemma~\ref{lem:cofinite-rational}, it
represents a member of \(\mathcal H_l\), so \(\chi(k)\) contains a
cofinite subset of \(l\).  An infinite proper subfield cannot be
cofinite: if \(a\notin\chi(k)\), the infinite coset \(a+\chi(k)\) is
disjoint from it.  Thus \(\chi(k)=l\).  Taking \(\sigma=\chi\), the
equality \(\sigma=hB=hMB_t\) outside a finite set proves
\eqref{eq:field-alignment} for our initial \(t\).

\emph{The same constant isomorphism works for every good function.}
For any other good function \(r\), Lemma~\ref{lem:common-partner}
supplies one good \(s_0\) that is algebraically independent of both
\(t\) and \(r\), with \(t+cs_0\) and \(r+cs_0\) good for every \(c\ne0\).
Lemma~\ref{lem:field-affine} gives
\(B_{s_0}B_t^{-1},B_{s_0}B_r^{-1}\in\mathcal H_l\), hence
\(B_rB_t^{-1}\in\mathcal H_l\).  The alignment already obtained for \(t\)
therefore gives \eqref{eq:field-alignment} for \(r\), with the same
\(\sigma\).
\end{proof}

\subsection{The isomorphism of perfect closures}

Fix \(\sigma\) and the \(H_t\) given by
Proposition~\ref{prop:constant-alignment}, and define
\[
 u_t=H_t^{-1}(t')\in L^i
 \qquad\text{for every good function }t.
\]
These are field elements, obtained by projective linear operations and
Frobenius powers.  They satisfy
\begin{equation}\label{eq:field-good-perfect}
 l(u_t)^i=l(t')^i.
\end{equation}
We show that the assignments \(t\mapsto u_t\) preserve every algebraic
relation.  This is the point where the finite exceptions in the individual
coordinate alignments cease to matter.

\begin{theorem}\label{thm:field-isomorphism}
There is a field isomorphism
\[
 \alpha:K^i\xrightarrow{\ \sim\ }L^i,\qquad \alpha|_k=\sigma,
\]
such that \(\alpha(t)=u_t\) for every good function \(t\).
In particular, if \(k(t)\) and \(l(t')\) are matched good rational
subfields, then \(\alpha(k(t)^i)=l(t')^i\).
\end{theorem}

\begin{proof}
Let \(\mathbf t=(t_1,\ldots,t_m)\) be a finite tuple of good functions,
and let \(Z_{\mathbf t}\subset(\Pj^1_l)^m\) be the relation variety of
\((u_{t_1},\ldots,u_{t_m})\).  We first prove
\begin{equation}\label{eq:field-relations}
 \sigma\bigl(X_{\mathbf t}(k)\bigr)=Z_{\mathbf t}(l),
\end{equation}
where \(\sigma\) acts coordinatewise.

The point map
\[
 \mathcal H=\prod_i H_{t_i}^{-1}:
 (\Pj^1(l))^m\longrightarrow(\Pj^1(l))^m
\]
is a Zariski homeomorphism.  Projective linear maps are isomorphisms; in
positive characteristic the coordinate map
\([X:Y]\mapsto[X^p:Y^p]\) is finite, radicial and surjective, hence a
homeomorphism on point spaces, as is its inverse.
Products with different powers in different coordinates have the same
property.  It follows that
\[
 \mathcal H\bigl(Y_{\mathbf t}(l)\bigr)=Z_{\mathbf t}(l).
\]
To justify this last equality at the level of relation varieties, place
the finite tuple of \(u_{t_i}\) in one finite purely inseparable extension
of \(L\).  On a common model, the tuples \((t_i')_i\) and
\((u_{t_i})_i\) are related by \(\mathcal H\) wherever they are defined.
Their images are dense in their respective relation varieties.
Taking closures under the homeomorphism gives the equality.

Apply now the first inclusion of
Lemma~\ref{lem:simultaneous-values}.  Delete from its source open the
finitely many exceptional values in each coordinate where
\eqref{eq:field-alignment} fails.  Each deleted coordinate fiber is a
proper closed subset, since every \(t_i\) is nonconstant.  The resulting
open is still dense, and on it
\(\mathcal H\circ\prod_iB_{t_i}=\sigma\).
Consequently a dense subset of \(\sigma(X_{\mathbf t}(k))\) lies in
\(Z_{\mathbf t}(l)\).  The former is the point set of the variety obtained
from \(X_{\mathbf t}\) by applying \(\sigma\) to coefficients, so taking
closures gives
\[
 \sigma(X_{\mathbf t}(k))\subset Z_{\mathbf t}(l).
\]

For the reverse inclusion, use the second inclusion of
Lemma~\ref{lem:simultaneous-values}.  Delete the finitely many target
coordinate values at which
\(B_{t_i}^{-1}=\sigma^{-1}H_{t_i}^{-1}\) fails.  Again this leaves a
dense open, now in \(Y_{\mathbf t}\), because each target coordinate is
nonconstant.  Its image under \(\mathcal H\) is dense in
\(Z_{\mathbf t}\) and lies in \(\sigma(X_{\mathbf t}(k))\).
Taking closures proves \eqref{eq:field-relations}.

Thus, for every polynomial \(P\in k[T_1,\ldots,T_m]\),
\[
 P(t_1,\ldots,t_m)=0
 \quad\Longleftrightarrow\quad
 \sigma(P)(u_{t_1},\ldots,u_{t_m})=0,
\]
where \(\sigma(P)\) is obtained by applying \(\sigma\) to the coefficients.
Good functions generate \(K\) by Lemma~\ref{lem:good-functions}.
The displayed equivalence for every finite tuple makes the assignments
\(a\mapsto\sigma(a)\) for \(a\in k\) and \(t\mapsto u_t\) for good \(t\)
well defined on all polynomial expressions, and preserves their nonzero
values.  Passing to quotients therefore gives an embedding
\(\alpha_0:K\hookrightarrow L^i\).

Because \(L^i\) is perfect, this embedding extends uniquely to an embedding
\(\alpha:K^i\hookrightarrow L^i\).  Its image contains \(l=\sigma(k)\)
and every \(u_t\), and is itself perfect; by
\eqref{eq:field-good-perfect} it contains the perfection of every matched
good rational subfield of \(L\).
All good rational subfields are matched, and their generators generate
\(L\).  The image of \(\alpha\) therefore contains \(L^i\), proving
surjectivity and the final assertion.
\end{proof}
